\documentclass[10pt,a4paper]{amsart}
\usepackage[margin=19mm]{geometry}
\usepackage{amsmath,amssymb,mathtools}
\usepackage[scr=rsfs,cal=euler]{mathalfa}
\usepackage{xfrac}
\usepackage[unicode,hidelinks,bookmarks=false]{hyperref}
\newcommand{\ie}{i.e.\ }
\newcommand{\cf}{cf.\ }
\newcommand{\resp}{respectively\ }
\newcommand{\Subsection}{\subsection}
\providecommand{\nobreakdash}{\nobreak\hbox{-}}
\setlength{\emergencystretch}{2em}
\multlinegap0pt
\usepackage{tcolorbox}
\usepackage{amssymb}
\usepackage{mathtools}
\usepackage{float}
\usepackage{xcolor}
\usepackage{bm}
\newtheorem{proposition}{Proposition}
\title{\bf superintegrable systems from Cartan commutants}
\author{Kai Jiang, Guorui Ma, Ian Marquette, Junze Zhang, Yao-Zhong Zhang}
\date{Source: arXiv:2605.14490v1 (CC BY 4.0)}
\begin{document}
\maketitle

\noindent\emph{Source and licence.} \bf superintegrable systems from Cartan commutants. Version arXiv:2605.14490v1 (CC BY 4.0), distributed by arXiv under the Creative Commons Attribution 4.0 International licence, http://creativecommons.org/licenses/by/4.0/, as recorded in the arXiv licence metadata for this exact version and shown on the abstract page https://arxiv.org/abs/2605.14490v1. Full licence notice: \texttt{licenses/arxiv-2605.14490-CC-BY-4.0.txt}.\par\smallskip

\noindent\textit{Editorial scope.} Counted unit: the article's proposition proposition (source0.tex lines 1479--1495). The article's complete original proof of it is delivered with it (source0.tex lines 1496--1525, one \texttt{proof} environment of 30 lines). No mathematics is written, repaired or re-derived: the statement and the proof are the article's own lines, in the article's own notation, and no retained line is substituted or re-flowed.  Retained uncounted as the source-local context the proof needs: lines 477--497.  Nothing was omitted from any retained span, so the package carries no omission record.  No retained line contains a citation, so the wrapper carries no bibliography and no bibliography entry.  No retained line contains a figure environment, an image-inclusion command or drawing code, so no image is delivered and the package declares no asset dependency.  Editorial formatting only, in addition: an AMS \texttt{amsart} preamble that restates the article's own declarations and macros used by the retained lines, namely the environments \texttt{proposition} on separate counters, together with the standard packages those lines need.  The retained environments print in this document as Proposition 1 in order of appearance, whereas the article prints the same items with the numbering of its own full document.  The complete native source of the article as distributed in the arXiv e-print for this exact version is bundled unchanged as \texttt{sources/200591/source0.tex}.
\begin{proposition}
\label{prop:equ}
Let $G$ be a $n$-dimensional semisimple Lie algebra $\mathfrak{g}$ with a Killing form $B$. %, and put $g:= B$. 
 Let $A\subset G$ be a connected closed subgroup with a Lie algebra $\mathfrak{a}\subset\mathfrak{g}$. Choose the basis of $\mathfrak{a}$ to be $\{X_1,\dots,X_s\}$ for some $s\leq n$. For any $f\in S(\mathfrak{g})$, define the adjoint action of $A$ on polynomials by
\begin{align*}
(a \cdot  f)(X):=f\big(aXa^{-1}\big),\qquad a\in A .
\end{align*}
Then for all $X\in\mathfrak{g},\ H\in\mathfrak{a}$ and $f \in S(\mathfrak{g})^A$, the following arguments are equivalent:
\begin{align*}
\text{\emph{(i)}}\quad & df_X \big([H,X]\big)=0\ ;\\
\text{\emph{(ii)}}\quad &
\sum_{i,k=1}^n  C_{jk}^i x_k \frac{\partial f}{\partial x_i} =0
\ \ \text{for each }j=1,\dots,s;\\
\text{\emph{(iii)}}\quad & B \big(\nabla f(X),[H,X]\big)=0\ ,
\end{align*} where $\nabla f$ is the gradient characterised by $df_X(V)=B(\nabla f(X),V)$, and for all $V\in\mathfrak{g}$ and $C_{jk}^i$ are structure constant. Equivalently, define
\begin{align*}
\mathcal{L}_j:=\sum_{i=1}^n\sum_{k=1}^n C_{jk}^i\,x_k\,\frac{\partial}{\partial x_i}\quad  j=1,\dots,s ,
\end{align*} we have \begin{align}
S(\mathfrak{g})^A=S(\mathfrak{g})^{\mathfrak{a}}=\bigcap_{j=1}^s \ker(\mathcal{L}_j). \label{eq:pdefin}
\end{align}
\end{proposition}

\begin{proposition}
    Let $S(\mathfrak{g})^T$ be the $T$-invariant polynomial Poisson subalgebra of $S(\mathfrak{g})$. Then \begin{align*}
        S(\mathfrak{g})^T \cong \mathbb{C}[h_1,\ldots,h_{n-1},p_{i_1,\ldots,i_d} : 2 \leq d \leq n] \big/ \mathcal{I} ,
\end{align*} where $\mathcal{I}$ is generated by the following relation \begin{align*}
    \prod_{u=1}^{k-1} p_{i_u,i_{u+1}}\, p_{i_k,i_1}
 &= p_{i_1,i_2,\ldots,i_k}\, p_{i_1,i_k,i_{k-1},\ldots,i_2},
 \notag\\
 \prod_{1 \leq i < j \leq n} p_{i,j}
 &= p_{1,2,\ldots,n}\, p_{1,n-1,\ldots,2}\,
 \prod_{m=1}^{n-2}\ \prod_{s=m+2}^{n} p_{m,s},
 \notag\\
 \prod_{i_1 \neq i_2 \neq \cdots \neq i_k} p_{i_1,i_2,\ldots,i_k}
 &= \left(\prod_{r<s} p_{r,s}\right)^{\varphi(k)},
 \qquad
 \varphi(k)=\prod_{s=2}^{k-1}(n-s). 
\end{align*}
\end{proposition}

\begin{proof}
  Since $T$ is connected, Proposition \ref{prop:equ} implies that $S(\mathfrak{g})^T = S(\mathfrak{g})^\mathfrak{t}$. Since $\mathfrak{t}$ is the Cartan subalgebra of diagonal traceless matrices. Write  \begin{align*}
      H=\mathrm{diag}(\lambda_1,\dots,\lambda_n)\in t, \qquad \lambda_1 +\cdots + \lambda_n = 0,
  \end{align*} such that the scalars $\lambda_i$ are the diagonal entries of $H$. Recall that, for an elementary matrices $e_{ij}$ ($i \neq j$), we have $[H,e_{ij}] = (\lambda_i - \lambda_j) e_{ij}$. Hence, by the PBW theorem, the monomial $p = \prod_{ i \neq j} e_{ij}^{m_{ij}} \in S(\mathfrak{g})^T$, then using the Leibniz rule, we have \begin{align*}
      \{H,p\}&=\sum_{i\neq j} m_{ij}\,\{H,e_{ij}\}\,e_{ij}^{m_{ij}-1}\prod_{(a,b)\neq(i,j)} e_{ab}^{m_{ab}}\\
&=\left(\sum_{i\neq j} m_{ij}(\lambda_i-\lambda_j)\right)p = 0.
  \end{align*} It is convenient to regroup the coefficient of $M$ as
\begin{align*}
\sum_{i\neq j} m_{ij}(\lambda_i-\lambda_j)
=\sum_{k=1}^n \lambda_k\Bigl(\underbrace{\sum_{j\neq k} m_{kj}}_{\text{out}(k)}-\underbrace{\sum_{i\neq k} m_{ik}}_{\text{in}(k)}\Bigr).
\end{align*}
Since $H$ varies over all diagonal traceless matrices, the scalars $\lambda_1,\dots,\lambda_n$ are arbitrary, subject only to $\sum_k\lambda_k =0$. Therefore, $p\in S(\mathfrak{g})^T$ if and only if  \begin{equation}
    \sum_{j\neq k}m_{kj}=\sum_{i\neq k}m_{ik} \label{eq:balanceequation}
\end{equation}   for every $k\in\{1,\dots,n\}$.

Define the root lattices by $R = \bigoplus_{i\neq j}\mathbb{Z}\,\varepsilon_{ij}$. Define also a $\mathbb{Z}$-linear map \begin{align*}
    \delta: R \rightarrow \mathbb{Z}^n, \quad \delta(\varepsilon_{ij}) = e_i  - e_j.
\end{align*} It is clear that $\delta$ is a group homomorphism.  Then \begin{align*}
\delta\Bigl(\sum_{i\neq j} m_{ij} \varepsilon_{ij}\Bigr) = 0
\end{align*} if and only if \eqref{eq:balanceequation} holds. For pairwise distinct indices $i_1,\dots,i_d$ ($d \geq 2$) define the cycle element by \begin{align*}
c_{i_1,\dots,i_d} := \varepsilon_{i_1 i_2} + \varepsilon_{i_2 i_3}+ \cdots+\varepsilon_{i_d i_1}\in R.
\end{align*} A direct telescoping computation gives $\delta(c_{i_1,\dots,i_d}) = 0$, so $c_{i_1,\dots,i_d}\in\ker(\delta)$.

We claim that any vector $$m := \sum_{i\neq j} m_{ij}\varepsilon_{ij}\in \ker(\delta)$$ with $m_{ij}\geq 0$ is a sum of such cycle elements by induction on $|m| := \sum_{i\neq j} m_{ij}$. If $|m| = 0$, there is nothing to prove. Otherwise, choose $i_1\neq i_2$ with $m_{i_1 i_2}>0$. Suppose $i_r$ has been chosen. The $i_r$-th balance equation implies that the total outgoing multiplicity from $i_r$ equals the total incoming multiplicity. Hence, in particular, there exists some $i_{r+1}\neq i_r$ with $m_{i_r i_{r+1}}>0$. Continuing in this way produces a sequence $i_1,i_2,i_3,\dots$. Since $\{1,\dots,n\}$ is finite, some index repeats: $i_s=i_t$ with $s<t$. If the repeated segment contains further repetitions, delete subloops to obtain a subsequence with pairwise distinct indices. In either case, we obtain a cycle $(j_1, \dots,j_d)$ with pairwise distinct entries such that each edge $(j_r \to j_{r+1})$ (with $j_{d+1} = j_1$) occurs with positive multiplicity in $m$. Therefore, $m - c_{j_1,\dots,j_d}$ still has nonnegative coefficients and still lies in $\ker (\delta)$ (since $\delta (c_{j_1,\dots, j_d})=0$). By induction, $m-c_{j_1,\dots,j_d}$ is a sum of cycle elements, hence so is $m$.

Translating back, subtracting $c_{j_1,\dots,j_d}$ corresponds to factoring $M$ by the cycle monomial $p_{j_1,\dots,j_d}$. Thus, any $T$-invariant monomial $M$ is a product of cycle monomials $p_{i_1,\dots,i_d}$ with pairwise distinct indices.

Since the action of $T$ on $t$ is trivial, every polynomial in $h_1,\dots,h_{n-1}$ is $T$-invariant. It follows that
$S(\mathfrak{g})^T$ is generated by $h_1,\dots,h_{n-1}$ together with the cycle monomials.
\end{proof}

\end{document}
